% --- Archon physics formalization source begin ---
% archon:physics
% archon:covers PhyXMiniProblems/problem_phyx_mini_0558.lean
% archon:source-report reports/phyx_mini/problem_phyx_mini_0558.source.json

\chapter{Physics problem phyx\_mini\_0558}
\label{ch:PhyXMiniProblems_problem_phyx_mini_0558}

\paragraph{Problem source.}
\#\# Physical scenario

Suppose that two cosmic-ray protons approach Earth from opposite directions as shown in figure. The speeds relative to Earth are measured to be \textbackslash{}( v\_1 = 0.6c \textbackslash{}) and \textbackslash{}( v\_2 = -0.8c \textbackslash{}).

\#\# Question

What is the velocity of each proton relative to the other?

\#\# Answer choices

- A: A: \textbackslash{}( 0.91c \textbackslash{})

- B: B: \textbackslash{}( 0.82c \textbackslash{})

- C: C: \textbackslash{}( 0.68c \textbackslash{})

- D: D: \textbackslash{}( 0.95c \textbackslash{})

\#\# Figure caption (auxiliary; use the image as primary evidence)

The image is a diagram showing two objects labeled "1" and "2" with velocity vectors indicated as \textbackslash{}(v\_1\textbackslash{}) and \textbackslash{}(v\_2\textbackslash{}), respectively. Both objects appear as circles. Object "1" is moving toward the left and object "2" toward the right. 

In the center, there is a larger blue circle labeled "Earth," representing the planet. A dashed line is drawn diagonally across this circle, possibly indicating an axis or trajectory. 

The arrows associated with \textbackslash{}(v\_1\textbackslash{}) and \textbackslash{}(v\_2\textbackslash{}) denote the direction of movement for each object. The text present in the image includes "1," "v\_1," "Earth," "v\_2," and "2."

\#\# Dataset metadata

Relativity; Multi-Formula Reasoning, Spatial Relation Reasoning

\paragraph{Recorded answer/context.}
D: D: \textbackslash{}( 0.95c \textbackslash{})

\paragraph{Figure/image path.}
/root/proposal\_for\_physic/science-mango/phyx\_mini\_run/phyx\_data/test\_image/558.png

\paragraph{Formalization target.}
create a compiling Lean file with sorry bodies at `PhyXMiniProblems/problem\_phyx\_mini\_0558.lean`. The Lean declarations must preserve the physical quantities, dimensions or dimensional roles, figure labels, governing-law hypotheses, and final relation expressed by this problem.
Use Mathlib/Physlib names found through LeanExplore where available. If a physics API is missing, introduce faithful local abstractions rather than scalar placeholder aliases.

\begin{theorem}[Physics formalization target]
\label{thm:physics:phyx_mini_0558:target}
\lean{PhyXMiniProblems.ProblemPhyXMini0558.problem_phyx_mini_0558}
\uses{def:physics:phyx-mini-0558:phyxminiproblems-problemphyxmini0558-opposingcosmicrayprotonssetup, def:physics:phyx-mini-0558:phyxminiproblems-problemphyxmini0558-matchesprimaryfigure, def:physics:phyx-mini-0558:phyxminiproblems-problemphyxmini0558-matchesproblemdata, def:physics:phyx-mini-0558:phyxminiproblems-problemphyxmini0558-hasphysicalinputparameters, def:physics:phyx-mini-0558:phyxminiproblems-problemphyxmini0558-satisfiesreciprocaleinsteinvelocitytransformations, def:physics:phyx-mini-0558:phyxminiproblems-problemphyxmini0558-relativevelocityfraction, def:physics:phyx-mini-0558:phyxminiproblems-problemphyxmini0558-isuniquematchinganswerchoice}
The assigned autoformalize agent should translate this physics problem into Lean declarations in the covered file, with theorem and lemma proof bodies written as `by sorry`.
\end{theorem}
\begin{proof}
This is an autoformalization task, not a proof task. Produce statements that can later be proved by the physics prover without weakening the physical model.
\end{proof}

% --- Archon declaration topology begin ---
\section{Declaration topology}
\begin{definition}[Signed Velocity Quantity]
  \label{def:physics:phyx-mini-0558:phyxminiproblems-problemphyxmini0558-signedvelocityquantity}
  \lean{PhyXMiniProblems.ProblemPhyXMini0558.SignedVelocityQuantity}
  A signed one-dimensional physical velocity component
\end{definition}
\begin{proof}
  This is the declared model definition of Signed Velocity Quantity.
\end{proof}
\begin{definition}[signed Velocity Readout]
  \label{def:physics:phyx-mini-0558:phyxminiproblems-problemphyxmini0558-signedvelocityreadout}
  \lean{PhyXMiniProblems.ProblemPhyXMini0558.signedVelocityReadout}
  \uses{def:physics:phyx-mini-0558:phyxminiproblems-problemphyxmini0558-signedvelocityquantity}
  Read a signed velocity component in selected length and time units
\end{definition}
\begin{proof}
  This is the declared model definition of signed Velocity Readout.
\end{proof}
\begin{definition}[signed Velocity In Meters Per Second]
  \label{def:physics:phyx-mini-0558:phyxminiproblems-problemphyxmini0558-signedvelocityinmeterspersecond}
  \lean{PhyXMiniProblems.ProblemPhyXMini0558.signedVelocityInMetersPerSecond}
  \uses{def:physics:phyx-mini-0558:phyxminiproblems-problemphyxmini0558-signedvelocityquantity, def:physics:phyx-mini-0558:phyxminiproblems-problemphyxmini0558-signedvelocityreadout}
  Meter-per-second readout of a signed physical velocity component
\end{definition}
\begin{proof}
  This is the declared model definition of signed Velocity In Meters Per Second.
\end{proof}
\begin{definition}[vacuum Light Speed In Meters Per Second]
  \label{def:physics:phyx-mini-0558:phyxminiproblems-problemphyxmini0558-vacuumlightspeedinmeterspersecond}
  \lean{PhyXMiniProblems.ProblemPhyXMini0558.vacuumLightSpeedInMetersPerSecond}
  \uses{def:physics:phyx-mini-0558:phyxminiproblems-problemphyxmini0558-signedvelocityinmeterspersecond}
  Physlib's exact dimensionful vacuum speed of light, read in SI units
\end{definition}
\begin{proof}
  This is the declared model definition of vacuum Light Speed In Meters Per Second.
\end{proof}
\begin{definition}[velocity Fraction Of Light]
  \label{def:physics:phyx-mini-0558:phyxminiproblems-problemphyxmini0558-velocityfractionoflight}
  \lean{PhyXMiniProblems.ProblemPhyXMini0558.velocityFractionOfLight}
  \uses{def:physics:phyx-mini-0558:phyxminiproblems-problemphyxmini0558-signedvelocityquantity, def:physics:phyx-mini-0558:phyxminiproblems-problemphyxmini0558-signedvelocityinmeterspersecond, def:physics:phyx-mini-0558:phyxminiproblems-problemphyxmini0558-vacuumlightspeedinmeterspersecond}
  A signed axial velocity expressed as a dimensionless fraction of c
\end{definition}
\begin{proof}
  This is the declared model definition of velocity Fraction Of Light.
\end{proof}
\begin{definition}[Figure Object]
  \label{def:physics:phyx-mini-0558:phyxminiproblems-problemphyxmini0558-figureobject}
  \lean{PhyXMiniProblems.ProblemPhyXMini0558.FigureObject}
  The three physical objects distinguished in the supplied bitmap
\end{definition}
\begin{proof}
  This is the declared model definition of Figure Object.
\end{proof}
\begin{definition}[Physical Object Kind]
  \label{def:physics:phyx-mini-0558:phyxminiproblems-problemphyxmini0558-physicalobjectkind}
  \lean{PhyXMiniProblems.ProblemPhyXMini0558.PhysicalObjectKind}
  The physical kind of an object in the scenario
\end{definition}
\begin{proof}
  This is the declared model definition of Physical Object Kind.
\end{proof}
\begin{definition}[Proton Label]
  \label{def:physics:phyx-mini-0558:phyxminiproblems-problemphyxmini0558-protonlabel}
  \lean{PhyXMiniProblems.ProblemPhyXMini0558.ProtonLabel}
  The two proton labels used by the prose and figure
\end{definition}
\begin{proof}
  This is the declared model definition of Proton Label.
\end{proof}
\begin{definition}[Inertial Frame Label]
  \label{def:physics:phyx-mini-0558:phyxminiproblems-problemphyxmini0558-inertialframelabel}
  \lean{PhyXMiniProblems.ProblemPhyXMini0558.InertialFrameLabel}
  Inertial frames needed to interpret all four velocity measurements
\end{definition}
\begin{proof}
  This is the declared model definition of Inertial Frame Label.
\end{proof}
\begin{definition}[Horizontal Direction]
  \label{def:physics:phyx-mini-0558:phyxminiproblems-problemphyxmini0558-horizontaldirection}
  \lean{PhyXMiniProblems.ProblemPhyXMini0558.HorizontalDirection}
  Horizontal direction relative to the left-to-right orientation of the image
\end{definition}
\begin{proof}
  This is the declared model definition of Horizontal Direction.
\end{proof}
\begin{definition}[Opposing Protons Figure]
  \label{def:physics:phyx-mini-0558:phyxminiproblems-problemphyxmini0558-opposingprotonsfigure}
  \lean{PhyXMiniProblems.ProblemPhyXMini0558.OpposingProtonsFigure}
  \uses{def:physics:phyx-mini-0558:phyxminiproblems-problemphyxmini0558-figureobject, def:physics:phyx-mini-0558:phyxminiproblems-problemphyxmini0558-protonlabel, def:physics:phyx-mini-0558:phyxminiproblems-problemphyxmini0558-horizontaldirection}
  Typed qualitative evidence from the primary image. The horizontal coordinates preserve only visual ordering and have no interpretation as physical distances or simultaneous spacetime positions
\end{definition}
\begin{proof}
  This is the declared model definition of Opposing Protons Figure.
\end{proof}
\begin{definition}[Opposing Cosmic Ray Protons Setup]
  \label{def:physics:phyx-mini-0558:phyxminiproblems-problemphyxmini0558-opposingcosmicrayprotonssetup}
  \lean{PhyXMiniProblems.ProblemPhyXMini0558.OpposingCosmicRayProtonsSetup}
  \uses{def:physics:phyx-mini-0558:phyxminiproblems-problemphyxmini0558-signedvelocityquantity, def:physics:phyx-mini-0558:phyxminiproblems-problemphyxmini0558-figureobject, def:physics:phyx-mini-0558:phyxminiproblems-problemphyxmini0558-physicalobjectkind, def:physics:phyx-mini-0558:phyxminiproblems-problemphyxmini0558-inertialframelabel, def:physics:phyx-mini-0558:phyxminiproblems-problemphyxmini0558-horizontaldirection, def:physics:phyx-mini-0558:phyxminiproblems-problemphyxmini0558-opposingprotonsfigure}
  The physical quantities of the problem. In particular, the two requested relative velocities are independent dimensionful fields; they are not defined from an answer choice or from a solved numerical formula
\end{definition}
\begin{proof}
  This is the declared model definition of Opposing Cosmic Ray Protons Setup.
\end{proof}
\begin{definition}[Matches Primary Figure]
  \label{def:physics:phyx-mini-0558:phyxminiproblems-problemphyxmini0558-matchesprimaryfigure}
  \lean{PhyXMiniProblems.ProblemPhyXMini0558.MatchesPrimaryFigure}
  \uses{def:physics:phyx-mini-0558:phyxminiproblems-problemphyxmini0558-opposingcosmicrayprotonssetup}
  Facts read from the primary bitmap: proton 1 is left of Earth and its v₁ arrow points right, while proton 2 is right of Earth and its v₂ arrow points left. The dashed mark across Earth and the visible text are retained, but the drawing contributes no quantitative position scale
\end{definition}
\begin{proof}
  This is the declared model definition of Matches Primary Figure.
\end{proof}
\begin{definition}[Matches Problem Data]
  \label{def:physics:phyx-mini-0558:phyxminiproblems-problemphyxmini0558-matchesproblemdata}
  \lean{PhyXMiniProblems.ProblemPhyXMini0558.MatchesProblemData}
  \uses{def:physics:phyx-mini-0558:phyxminiproblems-problemphyxmini0558-velocityfractionoflight, def:physics:phyx-mini-0558:phyxminiproblems-problemphyxmini0558-opposingcosmicrayprotonssetup}
  Frame assignments, object kinds, directions, and the two Earth-frame velocity readouts stated in the problem. Neither requested relative velocity is assigned a value here
\end{definition}
\begin{proof}
  This is the declared model definition of Matches Problem Data.
\end{proof}
\begin{definition}[Has Physical Input Parameters]
  \label{def:physics:phyx-mini-0558:phyxminiproblems-problemphyxmini0558-hasphysicalinputparameters}
  \lean{PhyXMiniProblems.ProblemPhyXMini0558.HasPhysicalInputParameters}
  \uses{def:physics:phyx-mini-0558:phyxminiproblems-problemphyxmini0558-vacuumlightspeedinmeterspersecond, def:physics:phyx-mini-0558:phyxminiproblems-problemphyxmini0558-velocityfractionoflight, def:physics:phyx-mini-0558:phyxminiproblems-problemphyxmini0558-opposingcosmicrayprotonssetup}
  Physical domain conditions on the supplied Earth-frame inputs. They select the ordinary subluminal branch but do not determine either requested relative velocity
\end{definition}
\begin{proof}
  This is the declared model definition of Has Physical Input Parameters.
\end{proof}
\begin{definition}[Satisfies Reciprocal Einstein Velocity Transformations]
  \label{def:physics:phyx-mini-0558:phyxminiproblems-problemphyxmini0558-satisfiesreciprocaleinsteinvelocitytransformations}
  \lean{PhyXMiniProblems.ProblemPhyXMini0558.SatisfiesReciprocalEinsteinVelocityTransformations}
  \uses{def:physics:phyx-mini-0558:phyxminiproblems-problemphyxmini0558-velocityfractionoflight, def:physics:phyx-mini-0558:phyxminiproblems-problemphyxmini0558-opposingcosmicrayprotonssetup}
  The one-dimensional Einstein velocity transformation is applied in both directions. If βa and βb are signed Earth-frame velocity fractions, then the velocity of a in b's rest frame is (βa - βb) / (1 - βa * βb). These are unsolved governing relations among physical quantities. They contain no exact result, rounded answer value, or answer-choice label
\end{definition}
\begin{proof}
  This is the declared model definition of Satisfies Reciprocal Einstein Velocity Transformations.
\end{proof}
\begin{definition}[relative Velocity Fraction]
  \label{def:physics:phyx-mini-0558:phyxminiproblems-problemphyxmini0558-relativevelocityfraction}
  \lean{PhyXMiniProblems.ProblemPhyXMini0558.relativeVelocityFraction}
  \uses{def:physics:phyx-mini-0558:phyxminiproblems-problemphyxmini0558-velocityfractionoflight, def:physics:phyx-mini-0558:phyxminiproblems-problemphyxmini0558-protonlabel, def:physics:phyx-mini-0558:phyxminiproblems-problemphyxmini0558-opposingcosmicrayprotonssetup}
  The signed requested velocity fraction for the named moving proton
\end{definition}
\begin{proof}
  This is the declared model definition of relative Velocity Fraction.
\end{proof}
\begin{definition}[Answer Choice]
  \label{def:physics:phyx-mini-0558:phyxminiproblems-problemphyxmini0558-answerchoice}
  \lean{PhyXMiniProblems.ProblemPhyXMini0558.AnswerChoice}
  Labels of the four relative-speed choices printed with the problem
\end{definition}
\begin{proof}
  This is the declared model definition of Answer Choice.
\end{proof}
\begin{definition}[displayed Relative Speed Fraction]
  \label{def:physics:phyx-mini-0558:phyxminiproblems-problemphyxmini0558-displayedrelativespeedfraction}
  \lean{PhyXMiniProblems.ProblemPhyXMini0558.displayedRelativeSpeedFraction}
  \uses{def:physics:phyx-mini-0558:phyxminiproblems-problemphyxmini0558-answerchoice}
  The positive coefficient of c printed beside each answer label
\end{definition}
\begin{proof}
  This is the declared model definition of displayed Relative Speed Fraction.
\end{proof}
\begin{definition}[recorded Dataset Answer]
  \label{def:physics:phyx-mini-0558:phyxminiproblems-problemphyxmini0558-recordeddatasetanswer}
  \lean{PhyXMiniProblems.ProblemPhyXMini0558.recordedDatasetAnswer}
  \uses{def:physics:phyx-mini-0558:phyxminiproblems-problemphyxmini0558-answerchoice}
  The answer label recorded by the source dataset; it is not a premise
\end{definition}
\begin{proof}
  This is the declared model definition of recorded Dataset Answer.
\end{proof}
\begin{definition}[Rounds To Nearest Hundredth]
  \label{def:physics:phyx-mini-0558:phyxminiproblems-problemphyxmini0558-roundstonearesthundredth}
  \lean{PhyXMiniProblems.ProblemPhyXMini0558.RoundsToNearestHundredth}
  Agreement with a positive speed coefficient printed to two decimal places
\end{definition}
\begin{proof}
  This is the declared model definition of Rounds To Nearest Hundredth.
\end{proof}
\begin{definition}[Matches Answer Choice]
  \label{def:physics:phyx-mini-0558:phyxminiproblems-problemphyxmini0558-matchesanswerchoice}
  \lean{PhyXMiniProblems.ProblemPhyXMini0558.MatchesAnswerChoice}
  \uses{def:physics:phyx-mini-0558:phyxminiproblems-problemphyxmini0558-protonlabel, def:physics:phyx-mini-0558:phyxminiproblems-problemphyxmini0558-opposingcosmicrayprotonssetup, def:physics:phyx-mini-0558:phyxminiproblems-problemphyxmini0558-relativevelocityfraction, def:physics:phyx-mini-0558:phyxminiproblems-problemphyxmini0558-answerchoice, def:physics:phyx-mini-0558:phyxminiproblems-problemphyxmini0558-displayedrelativespeedfraction, def:physics:phyx-mini-0558:phyxminiproblems-problemphyxmini0558-roundstonearesthundredth}
  A displayed choice matches only when it is the two-decimal rounding of the speed magnitude of each proton relative to the other
\end{definition}
\begin{proof}
  This is the declared model definition of Matches Answer Choice.
\end{proof}
\begin{definition}[Is Unique Matching Answer Choice]
  \label{def:physics:phyx-mini-0558:phyxminiproblems-problemphyxmini0558-isuniquematchinganswerchoice}
  \lean{PhyXMiniProblems.ProblemPhyXMini0558.IsUniqueMatchingAnswerChoice}
  \uses{def:physics:phyx-mini-0558:phyxminiproblems-problemphyxmini0558-opposingcosmicrayprotonssetup, def:physics:phyx-mini-0558:phyxminiproblems-problemphyxmini0558-answerchoice, def:physics:phyx-mini-0558:phyxminiproblems-problemphyxmini0558-matchesanswerchoice}
  The choice is the unique displayed match for both relative speeds
\end{definition}
\begin{proof}
  This is the declared model definition of Is Unique Matching Answer Choice.
\end{proof}
\begin{lemma}[velocity of each proton relative to other exact]
  \label{lem:physics:phyx-mini-0558:phyxminiproblems-problemphyxmini0558-velocity-of-each-proton-relative-to-other-exact}
  \lean{PhyXMiniProblems.ProblemPhyXMini0558.velocity_of_each_proton_relative_to_other_exact}
  \uses{def:physics:phyx-mini-0558:phyxminiproblems-problemphyxmini0558-opposingcosmicrayprotonssetup, def:physics:phyx-mini-0558:phyxminiproblems-problemphyxmini0558-matchesproblemdata, def:physics:phyx-mini-0558:phyxminiproblems-problemphyxmini0558-satisfiesreciprocaleinsteinvelocitytransformations, def:physics:phyx-mini-0558:phyxminiproblems-problemphyxmini0558-relativevelocityfraction}
  Substitution of β₁ = 3/5 and β₂ = -4/5 into the two Einstein transformations yields equal-and-opposite signed relative velocities: β₁|₂ = 35/37 and β₂|₁ = -35/37
\end{lemma}
\begin{proof}
  The conclusion follows from the governing relations and figure readouts named in the statement of velocity of each proton relative to other exact.
\end{proof}
% --- Archon declaration topology end ---
% --- Archon physics formalization source end ---
